% ==================================================
% プリアンブル設定ファイル
% パッケージ読み込み・基本設定
% ==================================================

% --------------------------------------------------
% 日本語フォント設定
% --------------------------------------------------

\usepackage[haranoaji]{luatexja-preset}

% --------------------------------------------------
% 欧文フォント設定（クロスプラットフォーム対応）
% --------------------------------------------------

\usepackage{fontspec}
\IfFontExistsTF{Times New Roman}{
  % Windows環境
  \setmainfont{Times New Roman}
  \setsansfont{Arial}
  \setmonofont{Courier New}
}{
  % Linux/macOS環境
  \IfFontExistsTF{Liberation Serif}{
    % Liberation Fonts（Times New Romanに近似）
    \setmainfont{Liberation Serif}
    \setsansfont{Liberation Sans}
    \setmonofont{Liberation Mono}
  }{
    % TeX Gyre Fonts（TeX Live標準フォント）
    \setmainfont{TeX Gyre Termes}                     % Times New Roman互換
    \setsansfont{TeX Gyre Heros}                      % Arial互換
    \setmonofont{TeX Gyre Cursor}                     % Courier New互換
  }
}

% --------------------------------------------------
% 表関連（biblatex内部で使用される可能性あり）
% --------------------------------------------------

\usepackage{array}
\usepackage{longtable}
\usepackage{tabularx}

% --------------------------------------------------
% ハイパーリンク・相互参照
% --------------------------------------------------

\usepackage{xurl}
\usepackage{hyperref}
\hypersetup{
  colorlinks=true,
  linkcolor=black,
  citecolor=black,
  urlcolor=blue,
  pdflang={ja-JP}
}

% --------------------------------------------------
% 業績一覧設定
% --------------------------------------------------

\usepackage{ifthen}
\usepackage[
  backend=biber, style=ieee, sorting=none,
  maxbibnames=99, giveninits=true, doi=true,
  url=true, eprint=true, date=iso, urldate=iso,
  seconds=true
]{biblatex}

% 日本語文献用の著者名フォーマット
% （姓名順，フルネーム，姓名スペースなし，andなし）
\DeclareNameFormat{japanese-names}{%
  \namepartfamily
  \ifdefvoid{\namepartgiven}{}{\namepartgiven}%
  \ifthenelse{\value{listcount}<\value{liststop}}
    {\addcomma\space}
    {}%
}

% 日本語文献用のフィールドフォーマット（強調なし）
\DeclareFieldFormat[article]{
  journaltitle:japanese}{#1}
\DeclareFieldFormat[inproceedings]{
  booktitle:japanese}{#1}
\DeclareFieldFormat[book]{title:japanese}{#1}

% 日本語会議論文で"in: "を削除
\renewbibmacro*{in:}{%
  \iffieldequalstr{langid}{japanese}
    {}
    {\printtext{\bibstring{in}\intitlepunct}}%
}

% 著者名と強調の切り替え
\AtEveryBibitem{%
  \iffieldequalstr{langid}{japanese}{%
    \DeclareNameAlias{author}{japanese-names}%
    \DeclareFieldAlias{journaltitle}{
      journaltitle:japanese}%
    \DeclareFieldAlias{booktitle}{
      booktitle:japanese}%
    \DeclareFieldAlias{title}{title:japanese}%
  }{}%
}

% --------------------------------------------------
% 目次設定
% --------------------------------------------------

\renewcommand{\contentsname}{目次}

% 目次に目次自身を含める
\usepackage{tocbibind}

% --------------------------------------------------
% マイクロタイポグラフィ
% --------------------------------------------------

\usepackage{microtype}
